\section{Agent 市场与大厂舰队}

前面复杂性风险讲的是 Manus 内部，本章看外部市场。访谈后半段，Peak 评价了 OpenAI、Anthropic、Google Gemini、xAI、Meta、Thinking Machines、Ilya 新公司等。他的观察有一个共同点：模型竞争没有终点，产品竞争也没有简单终点，但不同公司会沿着不同长板形成分化。

\begin{figure}[H]
\centering
\includegraphics[width=0.94\textwidth]{agent-market-map.png}
\caption{Agent 市场地图：Coding、垂直 Agent、C 端 Agent、生态配套和模型公司。}
\end{figure}

\begin{knowledgebox}{读图：Agent 市场不是一个赛道}
Coding Agent、垂直 Agent、C 端 Agent、支付/权限生态、模型公司都在推进。Manus 更接近 C 端和生产力工具之间的位置，既需要模型能力，也需要产品心智和真实工作流。
\end{knowledgebox}

\subsection{AI 不再只竞争用户时长}

Peak 提出一个很有意思的判断：移动互联网最终竞争用户总时长，而 Agent 可以在减少直接交互时长的同时创造更多价值。Agent 在后台持续工作，用户不需要一直盯着它。因此 AI 时代的稳态不一定像移动互联网那样由“注意力总量”决定。

\begin{knowledgebox}{Agent 与移动互联网的不同}
移动互联网产品常常争夺用户时间；Agent 产品可能争夺的是用户授权、任务入口、后台执行和结果可信度。它不一定需要用户一直在线。
\end{knowledgebox}

\subsection{本章小结}

Agent 市场会分化为多个层次。Manus 的机会，不只是模型能力，而是能否成为用户主力生产力入口，同时保持简单和独特。

